\documentclass[12pt,a4paper]{article}

% Encodage et langue
\usepackage[utf8]{inputenc}
\usepackage[T1]{fontenc}
\usepackage[french]{babel}

% Mise en page
\usepackage[a4paper,margin=2.5cm]{geometry}
\usepackage{lmodern}
\usepackage{enumitem}
\usepackage[hidelinks]{hyperref}

\begin{document}

\title{Journal de bord - Projet 2 : Jeu du tee-shirt}
\author{Zéïnab Aly Camara}
\date{Novembre 2026}

\maketitle
\tableofcontents
\newpage

\section{Introduction}

Dans le cadre du projet 2 de l'enseignement HAI716I, nous devons
concevoir une simulation multi-agents autour du jeu du tee-shirt.
L'objectif est d'étudier l'évolution des couleurs portées par les
agents au fil des rencontres et d'observer les conditions dans
lesquelles une population peut parvenir à un consensus.

Le travail est réalisé en binôme avec Neha. Nous avons réparti
les responsabilités en deux parties :
\begin{itemize}
    \item Neha est responsable du moteur de simulation, des
    expériences et des tests associés.
    \item Je suis responsable de l'implémentation des règles
    de mise à jour, de leurs tests et de mon journal de bord.
\end{itemize}

\section{Réflexion sur l'architecture}

Avant de commencer l'implémentation, nous avons pris le temps
de réfléchir à l'organisation du projet et aux échanges entre
le moteur de simulation et les règles de mise à jour.

Nous avons notamment choisi de définir un fichier commun,
\texttt{interface.py}, afin de disposer d'un contrat partagé
et de limiter les problèmes d'intégration entre nos deux
parties du travail.

Nous avons identifié trois éléments principaux à définir
dans ce fichier :

\begin{itemize}
    \item \texttt{Observation} : représente une rencontre entre
    deux agents et conserve les informations relatives à leurs
    couleurs ainsi que le gain obtenu.

    \item \texttt{RegleMiseAJour} : classe mère dont hériteront
    les différentes règles. Elle définit notamment la méthode
    \texttt{choisir\_couleur()}, chargée de déterminer la nouvelle
    couleur d'un agent.

    \item \texttt{ParametresSimulation} : regroupe les paramètres
    nécessaires à l'exécution d'une simulation, comme le nombre
    d'agents, le nombre maximal de tours et la graine aléatoire.
\end{itemize}

La classe \texttt{Agent} sera quant à elle définie dans
\texttt{jeu.py}, sous la responsabilité de Neha. Elle représentera
chaque participant à la simulation, avec notamment sa couleur,
sa mémoire et sa règle de mise à jour.

\subsection{Choix de conception}

Nous avons choisi une approche orientée objet pour les règles,
en nous appuyant sur le principe d'une classe mère commune.
Cette organisation reprend une approche déjà utilisée dans
un précédent projet, où chaque stratégie héritait d'une classe
mère et redéfinissait une méthode obligatoire.

Pour notre simulation, cette méthode est
\texttt{choisir\_couleur(memoire, couleur\_actuelle, rng)}.
Elle reçoit les observations précédentes de l'agent, sa couleur
actuelle et le générateur aléatoire fourni par le moteur.
Elle retourne ensuite la couleur qu'il devra porter.

Nous avons également retenu le principe d'un générateur
aléatoire partagé avec les règles. Cela permet de contrôler
les tirages aléatoires et de reproduire une simulation à partir
d'une même graine.

Enfin, nous avons identifié une contrainte importante pour
le moteur : les changements de couleur doivent être appliqués
simultanément. Les nouvelles couleurs de tous les agents doivent
donc être calculées à partir de leur état au début du tour,
avant que les modifications soient effectivement appliquées.

Cette séparation permet de distinguer clairement les
responsabilités : le moteur organise les rencontres et les
mises à jour, tandis que les règles déterminent le comportement
des agents.

\end{document}
